\section{A \textit{Reporter} komponens}\label{chp:reporter}

    A \textit{Reporter} szoftver komponens felelős a teszt keretrendszer kimenetét különböző formátumokra hozni.
    Kiemelendő, hogy ez a komponens egy külön parancssori futtatható a teszt keretrendszertől. Ezt az architekturális
    döntés a \emph{Unix filozófia} alapján hoztam meg, miszerint:

    \begin{foreigndisplaycquote}{american}{bst-journal}
        Make each program do one thing well. To do a new job, build afresh rather than complicate old programs by adding
        new features.
    \end{foreigndisplaycquote}

    A~\ref{sec:project-structure} fejezetben említett két kimeneti formátumot támogatja a komponens, amelyek rendre:

    \begin{itemize}
        \item Humán olvasható, \textit{HTML} jelentés.
        \item Harmadik féltől származó programok által olvasható \textit{JUnit XML}.
    \end{itemize}

    A kimeneti formátumok kiválasztásánál először csak a \textit{JUnit XML} lehetőséget vettem figyelembe. Azonban
    hamar rá kellett hogy jöjjek, hogy egy szoftver eszköz mint a \textit{SuperSet} teszt keretrendszer, csak akkor
    teremt értéket ha a tesztelők könnyen vizualizálni tudják a kimenetét. Ha a \textit{HTML} kimenet nem valósult
    volna meg, akkor a keretrendszer felhasználói kénytelenek lettek volna harmadik féltől származó szoftvereket
    alkalmazni. Vagy az \textit{SQLite} adatbázis megnyitásához, és lekérdezéséhez, vagy a \textit{JUnit XML}
    vizualizálásához.

    A két formátum kezelésére két külön-külön álló szoftver komponenst implementáltam, ezek a \textit{HTMLReporter} és
    a \textit{JXmlReporter} komponensek. A komponensek közel azonos adaton dolgoznak.

    Mindkét komponensnek szüksége van egy aggregált statisztikára ez a \textit{Stats} típus. Ez a típus gyakorlatilag
    egy összefoglalót ad az egész teszt futtatásról. Az összefoglalóban szerepel az összes konfigurációval kiválasztott
    teszteset száma, valamint hogy ezek közül mennyi teszteset volt sikeres, sikertelen, vagy kihagyott. Ezenkívül a
    futtatás idejét is tartalmazza másodpercben mérve. A \textbf{\ref{fig:stats-implementation}. ábrán} látható a
    \textit{Stats} típus.

    \begin{figure}[ht]
        \centering
        \begin{minted}{python}
@dataclass(slots=True)
class Stats(tuple[float, int, int, int, int, int]):
    time: float = 0
    passes: int = 0
    fails: int = 0
    errors: int = 0
    skipped: int = 0
    tests: int = 0
    \end{minted}
        \caption{\label{fig:stats-implementation} A \textit{Stats} típus implementációja}
    \end{figure}

    \pagebreak
    \subsection{A \textit{JXmlReporter} komponens}

        A \textit{JxmlReporter} komponens mind működésében mind implementációjában szerényebb mint a
        \textit{HTMLReporter} társa. Ez a tény gyakorlatilag a formátum egyszerűségéből adódik.

        A formátumot először az \emph{Apache Ant\textsuperscript{\texttrademark}}~\cite{ant} használta, és
        szabványosította, majd \textit{de-facto} szabvány lett az automata szoftvertesztek eredményeinek tárolására.

        A \textbf{\ref{fig:jxml-report}. ábrán} látható példa egy a \textit{SuperSet} teszt keretrendszer által
        generált \textit{JUnit XML-re}.

        \begin{figure}[ht]
            \centering
            \begin{minted}{xml}
<?xml version="1.0" encoding="UTF-8"?>
<testsuites name="SuperSet" tests="3" errors="0" failures="0" skipped="0" time="1.164">
    <testsuite name="C99LANG-ARMV8-O0" tests="1" errors="0" failures="0" skipped="0" time="0.291">
        <test_case name="2/1/1/2/t01" time="0.291" />
    </testsuite>
    <testsuite name="C99LIB-ARMV8-O0" tests="1" errors="0" failures="0" skipped="0" time="0.416">
        <test_case name="4/1/2/t33fsnofloat" time="0.416" />
    </testsuite>
    <testsuite name="LOOPS-ARMV8-O0" tests="1" errors="0" failures="0" skipped="0" time="0.457">
        <test_case name="optimize/s314/ts314intrestrict" time="0.457" />
    </testsuite>
</testsuites>
        \end{minted}
            \caption{\label{fig:jxml-report} Példa a \textit{JUnit XML} jelentésre}
        \end{figure}

    \subsection{A \textit{HTMLReporter} komponens}

        A könnyű áttekintés érdekében az eredmények kategorizálásra kerülnek. A \textit{HTMLReporter} komponens hét
        különböző kategóriát különböztet meg. Ezek a következők:

        \begin{itemize}
            \item \textit{SKIPPED}: kihagyott teszteset.
            \item \textit{CFAIL}: fordítási hibás teszteset.
            \item \textit{XCFAIL}: ismert fordítási hibás teszteset.
            \item \textit{CPASS}: fordítási sikeres teszteset.
            \item \textit{RFAIL}: futtatási hibás teszteset.
            \item \textit{XRFAIL}: ismert futtatási hibás teszteset.
            \item \textit{RPASS}: futtatási sikeres teszteset.
        \end{itemize}

        A \textit{HTMLReporter} komponens strukturáltan prezentálja az információt a felhasználó felé. A struktúra egy
        fába van rendezve.

        Az első gyökér elem, egy fejléc van, ami felsorolja a teljes futásra vonatkozó statisztikákat. Ahogy haladunk
        beljebb a listában, szintén minden elem egy statisztika fejléccel kezdődik, azonban az a fejléc csak az arra a
        lista elemre szűrt statisztikákat jeleníti meg.

        A lista elemek a könnyű vizuális keresés érdekében, szemantikusan színezettek, a kategóriák alapján. Erősségi
        sorrend szerint, pirossal ha sikertelen teszteset van a kategóriában, zölddel ha minden teszteset sikeres, és
        sárgával ha minden teszteset ki lett hagyva.

        A lista elemek ezután a következő sorrendben ágyazódnak egymásba:

        \begin{enumerate}
            \item A tesztcsomag neve.
            \item A \textit{Target} neve és a hozzátartozó konfiguráció.
            \item A hét kategória bontás.
            \item A kategórián belül az eredményeket mutató táblázat.
        \end{enumerate}

        A fa szerkezet leveleiben lévő táblázatok, a kategóriájuknak megfelelő oszlopokat tartalmaznak. Általánosan
        elmondható, hogy minden táblázatban van egy \emph{név} oszlop, hogy a teszteset azonosítható legyen.

        A sikeresen futott teszteseteknél a futtatás ideje látható. A kihagyott teszteseteknél a lista ami alapján a
        teszt ki lett hagyva. A sikertelen teszteseteknél a futtatás ideje, a program kimenete mind az
        \textit{STDOUT-ra} és \textit{STDERR-re}, valamint egy konfigurációs fájl amivel reprodukálni lehet a bukást. A
        mellékletekben található \textbf{\ref{fig:html-report-passes}.~\ref{fig:html-report-skips}.
        és~\ref{fig:html-report-errors}. ábrák} bemutatják a különböző táblázat formátumokat egy \textit{HTML} jelentés
        részletén.
